\section{Three simple cycles in every collar}\label{rw:lower-bound-section}

We construct three cycles in one field on the regular block
$\bar B=0$, $\bar m=1$. A variation at a reversible center gives a
three-root closing profile. We first fix the scale and then choose a
finite coefficient perturbation that preserves all three roots and
their actual return paths. Finally, a sequence of scales makes each
cycle converge to the entire hemicycle. This construction is independent
of the upper bound.

\subsection{The fixed-scale center and its actual saddle ratio}

For $0<\epsilon\ll1$ take the original coefficients
$B=0$, $m=\epsilon$ and $a=c=d=0$. The exact first integral is
\begin{equation}\label{rw:lower:center-energy}
 \mathscr H_\epsilon(x,y)=\frac{x^2}{2}+V_\epsilon(y),\qquad
 V_\epsilon(y)=(1+\epsilon)y-\frac{\epsilon y^2}{2}
                    -(1+\epsilon)\log(1+y),\quad y>-1.
\end{equation}
Indeed $V_\epsilon'=y(1-\epsilon y)/(1+y)$, so its Lie derivative
cancels $x\dot x=-xy+\epsilon xy^2$.
The potential decreases from infinity to zero on $(-1,0)$ and
increases from zero to a nondegenerate maximum at $y=1/\epsilon$.
Consequently the bounded well consists of periodic ovals with
$0<\mathscr H_\epsilon<H_s$, where
\begin{equation}\label{rw:lower:saddle-energy}
 H_s=\frac1{2\epsilon}+1
                  -(1+\epsilon)\log(1+1/\epsilon).
\end{equation}
Its boundary is the homoclinic loop through $(0,1/\epsilon)$,
as in Proposition~\ref{rw:centers:reversible}.
The two loop intersections with $y=0$ have normalized coordinates
$P_s=Q_s=\sqrt{2\epsilon H_s}\to1$.

We keep $\bar B=0$, $\bar m=1$ and use
$(\bar a,\kappa,\nu)$ as local coefficient variables, where
$\kappa=r-1$ and $r$ is the actual upper saddle ratio.
For the finite-order constructions of Section~\ref{rw:transitions},
reserve at least six total input, coefficient and Euler scale orders
before choosing the common neighborhoods. The variation below uses
ordinary coefficient derivatives at a fixed positive scale.
The exact upper field before division by $Y+\epsilon$ is
\[
 X_\xi=Y(Y-1)+\epsilon\bar aX,\qquad
 Y_\xi=(Y+\epsilon)(X+\bar cY)+\nu Y^2.
\]
At the center the saddle is $(0,1)$ and its Jacobian is
$\left(\begin{smallmatrix}0&1\\1+\epsilon&0\end{smallmatrix}\right)$.
Let $S$ be the moving saddle and $T$ the trace of its Jacobian.
Differentiating both equilibrium equations gives
\[
 S_{\bar a}=(0,0),\qquad S_{\bar c}=(-1,0),\qquad
 S_\nu=(-1/(1+\epsilon),0),
\]
and, including this motion in the trace,
\[
 T_{\bar a}=\epsilon,\qquad T_{\bar c}=1+\epsilon,\qquad
 T_\nu=\frac{1+2\epsilon}{1+\epsilon}.
\]
Since $\kappa=-T/\lambda_u$ and $\lambda_u=\sqrt{1+\epsilon}$
at the center, the parameter inverse
$\bar c=\bar c_\epsilon(\bar a,\kappa,\nu)$ satisfies exactly
\begin{equation}\label{rw:lower:actual-ratio}
 \bar c_{\bar a}=-\frac{\epsilon}{1+\epsilon},\qquad
 \bar c_\kappa=-\frac1{\sqrt{1+\epsilon}},\qquad
 \bar c_\nu=-\frac{1+2\epsilon}{(1+\epsilon)^2}.
\end{equation}
The equilibrium and parameter Jacobians have uniform nonzero
margins on a fixed positive-height box. The implicit function
theorem therefore supplies this inverse with the finite coefficient
jets used below on a common neighborhood.

Choose the symmetric local saddle sections and coordinates of
Section~\ref{rw:upper-resonance}. Its passage is
\[
 v=t^{1+\kappa}e^{w(t)},\qquad t>0,
\]
and at the reversible center it is exactly $v=t$.
The two regular outer legs reach $Y=0$ with positive labels $p(t)$
and $q(v)$ near one. Increasing $t$ or $v$ moves into the period
annulus, giving $p_t,q_v<-c<0$ on smaller fixed intervals.

\begin{lemma}[The regular lower block]\label{rw:lower:regular-block}
With one fixed weak reference on each side, the two complete central
maps have the form $s_++U(t)$ and $s_-+V(v)$, where
\begin{equation}\label{rw:lower:regular-amplitudes}
 U(t)=e^{-S_+(t)/\epsilon+\ell_+(t)},\qquad
 V(v)=e^{-S_-(v)/\epsilon+\ell_-(v)}.
\end{equation}
On the indicated coefficient block the functions $S_\pm,\ell_\pm$
have bounded prescribed finite ordinary input, coefficient and Euler
scale jets. Moreover $S_\pm>s_*>0$ and
$-S_+',-S_-'>c>0$. The reference values are independent of input,
and
\begin{equation}\label{rw:lower:reference-gap}
 s_--s_+=\sqrt\epsilon\,\nu H,\qquad H>h_*>0,\qquad
 H|_{\bar a=\kappa=\nu=0}\longrightarrow\sqrt{2\pi}.
\end{equation}
Thus $U',V'>0$, and the rates
\[
 \mathsf A_+=\epsilon U''/U',\qquad
 \mathsf A_-=\epsilon V''/V'
\]
have positive margins and bounded finite mixed jets.
For each of the finitely many coefficient derivatives used here,
\begin{equation}\label{rw:lower:absolute-flat}
 |\partial_{\mathbf b}^{\alpha}U(0)|
 +|\partial_{\mathbf b}^{\alpha}V(0)|+u_\epsilon'(0)
 \le C\epsilon^{-M}e^{-c/\epsilon},
 \qquad u_\epsilon(t)=U(t;0)=V(t;0).
\end{equation}
At the center set
$A_\epsilon=\epsilon u_\epsilon''(0)/u_\epsilon'(0)$;
then $0<c\le A_\epsilon\le C$.
\end{lemma}

\begin{proof}
Both axis labels lie in $(3/4,5/4)$, and $1\pm\bar aX$ have
positive margins on a larger fixed interval. On the regular layer
from $z=1$ to a fixed small $z=d_0>0$, the normal speed has a
nonzero margin. The original graph equation with independent variable
$z$ gives $P=p+\epsilon R(p,\mathbf b,\epsilon)$ with bounded
finite jets, and likewise on the reflected side.

Choose one weak graph and keep it fixed as $p$ varies.
Proposition~\ref{rw:phase}
and $w=\epsilon zT$ give its height as $\epsilon\nu$ times a
function with bounded ordinary endpoint and coefficient jets on this
compact positive $X$ interval. Its restriction at each $P$ has the
same central reference. Apply Proposition~\ref{rw:transitions:central-core}
with $P$ as an independent regular endpoint: $X=Py$ puts its weighted
reference equation on a fixed interval, and $P$ has a positive lower
margin. Its relative amplitude and the positive input deviation
$d_0-z_b(P)$ therefore have bounded ordinary $P$ jets.
With
\[
 A_\pm(p)=\int_0^p\frac{X}{1\pm\bar aX}\,\dd X,
\]
the outgoing deviation is $a_+(p)e^{-A_+(p)/\epsilon}$, with
bounded mixed jets of $\log a_+$, since the layer correction
\[
 \frac{A_+(P)-A_+(p)}{\epsilon}
   =R(p,\mathbf b,\epsilon)
       \int_0^1 A_+'(p+s\epsilon R)\,\dd s
\]
has those bounds. Reflection gives the incoming statement. Composing
with the two regular upper legs gives \eqref{rw:lower:regular-amplitudes};
their jets and inverse margins are supplied by the low-layer estimate
of Section~\ref{rw:upper-resonance}, whose scale carrier has integral
$\epsilon\log(1+\delta/\epsilon)$. These connections exclude the
saddle transit. The signs $A_\pm'>c$ and $p_t,q_v<-c$ give the
action derivative margins.

For the reference difference, both references vanish when $\nu=0$.
Their Gaussian variational equations give a factor $\nu$ directly.
At the center, the first $\nu$ variation $\eta$ on the outgoing
core satisfies
\[
 \eta_X=\frac{X}{\epsilon(1+\epsilon)}\eta+\frac1{1+\epsilon}.
\]
A fixed positive endpoint contributes only an exponentially small
terminal term at $X=0$. Thus
\[
 \eta(0)=-\frac1{1+\epsilon}
       \int_0^{R_0}e^{-X^2/(2\epsilon(1+\epsilon))}\,\dd X
       +O(\epsilon^{-M}e^{-c/\epsilon}).
\]
The reflected reference has the opposite derivative. The change
$X=\sqrt\epsilon\,u$, with its Gaussian integrable majorant, proves
the limit in \eqref{rw:lower:reference-gap}.
On the larger small coefficient box the same kernel argument for
the exact difference quotient in $\nu$ retains this leading positive
term, with an error $O(\sqrt\epsilon)$ after division by
$\sqrt\epsilon$. More explicitly, put $X=\sqrt\epsilon\,u$ and
$z_b=\sqrt\epsilon\,\nu\zeta$ in
\eqref{rw:relative:core-field}, using the continuous quotient at
$\nu=0$. On each bounded $u$ interval its equation tends uniformly
to $\zeta_u=u\zeta+1$; the weighted reference estimate and Gaussian
kernel control the remaining tail. The outgoing value tends to
$-\sqrt{\pi/2}$ and the reflected value to $\sqrt{\pi/2}$,
uniformly in this coefficient box. This gives the uniform positive
margin for $H$.
These estimates use the weighted reference equation and its kernel,
without requiring ordinary high-order $X$ derivatives at the Gaussian
center.

Finally put $Z_\pm=-S_\pm'+\epsilon\ell_\pm'>c$.
Direct differentiation of \eqref{rw:lower:regular-amplitudes} gives
\[
 \epsilon U''/U'=Z_++\epsilon Z_+' /Z_+,
\]
and the corresponding identity for $V$. All required rate jets are
therefore bounded and positive after reducing the scale.
Each fixed coefficient derivative of the amplitudes costs only a
finite inverse power of $\epsilon$, multiplying $e^{-s_*/\epsilon}$.
This proves \eqref{rw:lower:absolute-flat} and the center rate bound.
\end{proof}

\subsection{A first variation with three simple closing zeros}

The actual closing difference in these coordinates is
\begin{equation}\label{rw:lower:closing}
 F(t)=\mathscr D(t)-\sqrt\epsilon\,\nu H,\qquad
 \mathscr D(t)=U(t)-V(v(t)).
\end{equation}
Define
\[
 K(t)=\epsilon\log\frac{V'(v(t))v'(t)}{U'(t)},\qquad D=t\partial_t.
\]
Integration of the lower logarithmic rates gives the exact formula
\begin{equation}\label{rw:lower:exact-logarithm}
\begin{aligned}
 K(t)&=K_c+\epsilon\kappa\log t+\mathcal E(t),\\
 K_c&=\epsilon\log\frac{(1+\kappa)V'(0)}{U'(0)},\\
 \mathcal E(t)&=\int_0^{v(t)}\mathsf A_-(s)\,\dd s
       -\int_0^t\mathsf A_+(s)\,\dd s
       +\epsilon\left[w+\log\frac{1+\kappa+Dw}{1+\kappa}\right].
\end{aligned}
\end{equation}
The constant term $K_c$ has the selected bounded coefficient jets.
Write, at the center and in the actual parameters,
\[
 a_c(\epsilon)=\partial_{\bar a}K_c,\qquad
 b_c(\epsilon)=\partial_\kappa K_c.
\]
Their limits are
\begin{equation}\label{rw:lower:constant-columns}
 a_c\longrightarrow-2/3,\qquad b_c\longrightarrow1.
\end{equation}
Indeed Lemma~\ref{rw:lower:regular-block} gives
$K_c=A_+(p_c)-A_-(q_c)+O_{C_{\mathbf b}^2}
(\epsilon(1+|\log\epsilon|))$ on $\nu=0$.
At zero scale,
\[
 p_c=\frac{\sqrt{\bar c^2+4}-\bar c}{2},\qquad
 q_c=\frac{\sqrt{\bar c^2+4}+\bar c}{2},
\]
and
\[
 A_+(p)=p^2/2-\bar ap^3/3+O(\bar a^2),\qquad
 A_-(q)=q^2/2+\bar aq^3/3+O(\bar a^2).
\]
Use \eqref{rw:lower:actual-ratio} to differentiate these expressions
at fixed $\kappa$ or fixed $\bar a$, respectively.
The first derivative is $-2\int_0^1X^2\,\dd X=-2/3$.
For the second,
$p_{c,\bar c}=-1/2$, $q_{c,\bar c}=1/2$ and
$\bar c_\kappa\to-1$, giving one.

We also record directly the upper first variation needed here.
In the prepared upper rectangle the orbital equation is
\[
 V_U=-rV/U-\sigma V^2-UV^3\mathcal R(U,V).
\]
At the center $V=t/U$, $r=1$ and $\sigma=\mathcal R=0$.
For any bounded coefficient direction, the derivatives
$\dot\sigma$ and $\dot{\mathcal R}$ are $O(\epsilon)$ with their
needed finite state jets, by the original-field construction in
Section~\ref{rw:upper-resonance}. Set $\dot V=(t/U)\eta(U)$.
Linearization of this equation, with its fixed input, gives
\[
 \eta_U=-\dot r/U-\dot\sigma\,t/U
                     -t^2\dot{\mathcal R}(U,t/U)/U,\qquad \eta(t)=0.
\]
At $U=1$, after removing $\dot r\log t$, this is
\begin{equation}\label{rw:lower:upper-linearization}
 \dot w(t)=\dot\sigma\,t\log t
     -t^2\int_t^1\dot{\mathcal R}(u,t/u)\,\frac{\dd u}{u},
 \qquad |D^j\dot w|\le C\epsilon t(1+|\log t|)\quad(0\le j\le3).
\end{equation}
Euler differentiation keeps $u,t/u$ bounded and adds bounded endpoint
terms, proving the displayed estimates for the true saddle variation.

The ratio variation supplies the logarithmic term in $K$. We choose
the $\bar a$ direction to balance its constant term at a smaller input
scale, and then use $\nu$ to shift the closing difference. These three
choices produce the following profile.

\begin{proposition}[The closing profile]\label{rw:lower-variation}
For every sufficiently small $\epsilon>0$, there is a bounded
direction $h_\epsilon\to3/2$, an exponentially small
$\beta_\epsilon$, and positive numbers $\delta_\epsilon,S_\epsilon$
such that along the original coefficient curve determined by
\[
 (\bar a,\kappa,\nu)=(h_\epsilon\theta,\theta,\beta_\epsilon\theta)
\]
the first variation $f_\epsilon=\partial_\theta F|_{\theta=0}$ obeys
\begin{equation}\label{rw:lower:profile-limit}
 \frac{f_\epsilon(\delta_\epsilon x)}{S_\epsilon}
 \longrightarrow
 P(x)=\frac{x^2}{2}-x\log x-x-\frac1{20}.
\end{equation}
The convergence is uniform on $[0,10]$ after continuous extension
of the values at zero, and is $C^1$ on every compact subinterval
of $(0,10]$. The curve uses the actual ratio inverse
\eqref{rw:lower:actual-ratio}. This proposition requires no bound on
the number of cycles of other fields.
\end{proposition}

\begin{proof}
Put $L=\log(1/\epsilon)$ and choose
\begin{equation}\label{rw:lower:direction-and-scale}
 h_\epsilon=
 \frac{\epsilon(L+\log L+\log A_\epsilon+2)-b_c}{a_c},\qquad
 \delta_\epsilon=\frac{\epsilon}{A_\epsilon L},\qquad
 S_\epsilon=u_\epsilon'(0)\delta_\epsilon>0.
\end{equation}
The limits \eqref{rw:lower:constant-columns} prove the assertion
about $h_\epsilon$. First vary
$(\bar a,\kappa,\nu)=(h_\epsilon\theta,\theta,0)$, and write
$k_\epsilon=\dot K$ and $d_\epsilon=\dot{\mathscr D}$.
At the center let $\mathsf A$ be the common lower rate;
$\mathsf A(t)=A_\epsilon+O(t)$ with bounded finite derivatives.
The variation of the endpoint is
$\dot v=t(\log t+\dot w)$. Differentiating every term in
\eqref{rw:lower:exact-logarithm} gives
\[
 \dot{\mathcal E}(t)
 =\int_0^t(\dot{\mathsf A}_--\dot{\mathsf A}_+)(s)\,\dd s
    +\mathsf A(t)t(\log t+\dot w)
    +\epsilon(\dot w+D\dot w).
\]
The explicit $\kappa$ derivatives in the logarithm's numerator and
denominator cancel, leaving the displayed moving-endpoint term.
The mixed rate bounds and \eqref{rw:lower:upper-linearization} imply
\begin{equation}\label{rw:lower:critical-error}
\begin{aligned}
 k_\epsilon(t)&=\epsilon(L+\log L+\log A_\epsilon+2)
                 +\epsilon\log t+A_\epsilon t\log t+R_\epsilon(t),\\
 |R_\epsilon(t)|&\le C\{t+t^2|\log t|
                  +\epsilon^2t(1+|\log t|)\},\\
 |R_\epsilon'(t)|&\le C\{1+t(1+|\log t|)
                  +\epsilon^2(1+|\log t|)\}.
\end{aligned}
\end{equation}
Here the rate-difference integral contributes $O(t)$, the error
$\mathsf A(t)-A_\epsilon$ contributes $O(t^2|\log t|)$, and
the upper relative terms give the last term in each bound.

The identity $\log\delta_\epsilon=-L-\log L-\log A_\epsilon$
now gives exactly
\[
 \frac{k_\epsilon(\delta_\epsilon x)}{\epsilon}
 -(2+\log x-x)
 =\frac{x}{L}(\log x-\log L-\log A_\epsilon)
       +\frac{R_\epsilon(\delta_\epsilon x)}{\epsilon}.
\]
It proves $C^1$ convergence to $G(x)=2+\log x-x$ on compact positive
intervals. The value error also tends to zero in $L^1(0,10)$:
\eqref{rw:lower:critical-error} supplies an integrable majorant
$C(1+|\log x|+x|\log x|)$ there, and quantitative integration gives
an error $O((1+\log L)/L)$.

The exact identity
$\mathscr D'=U'(1-e^{K/\epsilon})$, differentiated at the center,
is
\begin{equation}\label{rw:lower:derivative-identity}
 d_\epsilon'(t)=-u_\epsilon'(t)k_\epsilon(t)/\epsilon.
\end{equation}
Furthermore
\[
 \log\frac{u_\epsilon'(\delta_\epsilon x)}{u_\epsilon'(0)}
 =\frac1\epsilon\int_0^{\delta_\epsilon x}\mathsf A(s)\,\dd s
 =\frac{x}{L}+O(\epsilon/L^2),\qquad 0\le x\le10.
\]
The value $d_\epsilon(0)$ exists: the regular lower maps at the
separatrix inputs are smooth in coefficients, and $t\log t\to0$.
Integrating \eqref{rw:lower:derivative-identity} from zero gives
\begin{equation}\label{rw:lower:integrated-profile}
 \frac{d_\epsilon(\delta_\epsilon x)-d_\epsilon(0)}{S_\epsilon}
 \longrightarrow J(x)=\frac{x^2}{2}-x\log x-x.
\end{equation}
The integrable majorant proves uniform convergence of values on
$[0,10]$; the same identity proves $C^1$ convergence on compact
positive intervals, with error $O((1+\log L)/L)$.

It remains to realize the constant shift by an actual line-breaking
parameter. At fixed $(\bar a,\kappa)$ put
$e_\epsilon(t)=\partial_\nu\mathscr D(t)|_0$ and $H_\epsilon=H|_0$,
and define
\begin{equation}\label{rw:lower:line-breaking-direction}
 \beta_\epsilon=
  \frac{d_\epsilon(0)+S_\epsilon/20}
       {\sqrt\epsilon\,H_\epsilon-e_\epsilon(0)}.
\end{equation}
Lemma~\ref{rw:lower:regular-block} implies
\[
 \sqrt\epsilon\,H_\epsilon-e_\epsilon(0)
       \ge (h_*/2)\sqrt\epsilon,\qquad
 |\beta_\epsilon|\le C\epsilon^{-M}e^{-c/\epsilon}.
\]
On the curve in the statement, the first variation satisfies exactly
\[
 \frac{f_\epsilon(t)}{S_\epsilon}
 =\frac{d_\epsilon(t)-d_\epsilon(0)}{S_\epsilon}-\frac1{20}
       +\beta_\epsilon
          \frac{e_\epsilon(t)-e_\epsilon(0)}{S_\epsilon}.
\]
An absolute bound for the last difference would lose the exponentially
small normalization $S_\epsilon$. Instead we use the differentiated
closing identity. At fixed $\kappa$ the $\nu$ derivative of
\eqref{rw:lower:exact-logarithm} has no $\epsilon\log t$ term.
Its constant, rate-integral and upper terms are bounded, so
$|K_\nu(t)|\le C$ for $0<t\le10\delta_\epsilon$.
The differentiated exact identity gives
\[
 e_\epsilon'=-u_\epsilon'K_\nu/\epsilon,\qquad
 \left\|\frac{e_\epsilon(\delta_\epsilon x)-e_\epsilon(0)}
                    {S_\epsilon}\right\|_{C^1([a,b])}
       \le C_{a,b}/\epsilon
 \quad(0<a<b\le10).
\]
The value bound also holds on $[0,10]$. The derivative identity
cancels the common factor $u_\epsilon'(0)$; the exponentially small
$\beta_\epsilon$ then sends both normalized errors to zero.
Since $\nu=0$ at the center, differentiation of $\nu H$ introduces
only $\beta_\epsilon H_\epsilon$.
Combining these facts with \eqref{rw:lower:integrated-profile}
proves \eqref{rw:lower:profile-limit}; on each $[a,b]$ its error is
at most
\[
 C_{a,b}\left(\frac{1+\log L}{L}
                     +\frac{|\beta_\epsilon|}{\epsilon}\right).
\]
\end{proof}

\subsection{Finite perturbations and localization}

\begin{theorem}\label{rw:three-simple-cycles}
For every open neighborhood $\mathscr U$ of the hemicycle $\Gamma$
on the northern Poincare sphere, there is a sequence of quadratic
fields tending to the base field, each with three distinct simple
limit cycles whose northern images lie in $\mathscr U$.
Each of the three sequences of cycles can be chosen to converge to
the entire $\Gamma$ in Hausdorff distance.
Consequently the quadratic cyclicity of $\Gamma$ is at least three.
\end{theorem}

\begin{proof}
The profile has signs $-,+,-,+$ at $10^{-3},10^{-1},1,10$, and
\[
 P'(x)=x-\log x-2,\qquad P''(x)=1-1/x.
\]
Thus $P'$ has exactly two simple zeros, one on each side of one.
Rolle's theorem and the four signs give exactly three simple positive
zeros. Choose disjoint compact root intervals with $P'$ bounded
away from zero and opposite endpoint signs.

First choose $\epsilon$ so that Proposition~\ref{rw:lower-variation}
preserves these margins. At this fixed scale, smooth coefficient
dependence at positive section inputs gives
\[
 \frac{F(\delta_\epsilon x;\theta)}
                 {\theta S_\epsilon}
 \longrightarrow
 \frac{f_\epsilon(\delta_\epsilon x)}{S_\epsilon}
 \quad\hbox{in }C^1
\]
on the three intervals as $\theta\to0$, preserving one simple
closing zero in each for every sufficiently small positive $\theta$.

At the fixed scale, the center ovals with
$10^{-3}\delta_\epsilon\le t\le10\delta_\epsilon$ form a compact
family inside the bounded period annulus, with finite flights,
transverse section crossings and a positive minimum central height.
Continuous dependence preserves these properties for small $\theta$,
so all three closing zeros compare actual positive outputs for the
same field. The forward first return on the incoming saddle section is
\[
 \mathcal R(t)=U^{-1}\bigl(s_--s_++V(v(t))\bigr).
\]
Transversality preserves the center's single-turn itinerary.
At a closing zero,
\begin{equation}\label{rw:lower:actual-return-multiplier}
 \mathcal R'(t)=\frac{V'(v(t))v'(t)}{U'(t)}=e^{K(t)/\epsilon},
 \qquad F'(t)=U'(t)\bigl(1-\mathcal R'(t)\bigr).
\end{equation}
Since $U'>0$ and $F'\ne0$, the three zeros give three distinct
hyperbolic cycles, each of multiplicity one. The admissible $\theta$
may depend on $\epsilon$.

We next locate the center ovals before making the final perturbation.
At $\bar B=0$, $\bar m=1$, the saddle first integral used in the
upper coordinates is
\[
 \mathcal H=X^2-\int_1^Y\frac{2s(s-1)}{s+\epsilon}\,\dd s
       =2\epsilon(\mathscr H_\epsilon-H_s).
\]
The symmetric product coordinates satisfy $\mathcal H=-4UV$.
On the normalized local sections,
$UV=\chi_{\rm loc}^2t$, where the fixed section factor
$\chi_{\rm loc}$ and its reciprocal are bounded.
Thus, uniformly for the selected ovals,
\begin{equation}\label{rw:lower:energy-localization}
 H_s-\mathscr H_\epsilon
          =\frac{2\chi_{\rm loc}^2t}{\epsilon},
 \qquad \epsilon\mathscr H_\epsilon\longrightarrow\frac12.
\end{equation}
At the central section $x=0$, $z=s$, the identity
\[
 \mathscr H_\epsilon(0,s-1)
     =(1+\epsilon)(s-1-\log s)
                          -\epsilon(s-1)^2/2
\]
also shows $\epsilon(s-1-\log s)\to1/2$.

Write $C_{\epsilon,t}$ for the center oval with section input $t$.
We prove
\begin{equation}\label{rw:lower:hausdorff}
 \sup_{10^{-3}\delta_\epsilon\le t\le10\delta_\epsilon}
 d_H\bigl(\iota(C_{\epsilon,t}),\Gamma\bigr)\longrightarrow0,
\end{equation}
where $d_H$ is spherical Hausdorff distance. We prove both distance
estimates, so that the limit is the entire hemicycle.
First take any sequence of points on selected ovals
with $\epsilon\to0$. A bounded affine subsequence with a limit in
$y>-1$ is impossible: \eqref{rw:lower:center-energy} stays bounded
on compact subsets there, whereas the energies diverge uniformly.
A bounded affine limit must therefore lie on $y=-1$.
An unbounded sequence has $y>-1$, so every spherical accumulation
point lies on the upper equatorial semicircle. Both types lie in
$\Gamma$. Compactness proves that the maximum distance from the
selected ovals to $\Gamma$ tends to zero.

Conversely, let $E$ range over the selected energies, so that
$\epsilon E\to1/2$ uniformly. For every fixed $M$, the lower branch
of $V_\epsilon(y)=E-x^2/2$, $|x|\le M$, exists in $-1<y<0$.
The potential is uniformly bounded on compact subintervals of
$(-1,0]$, whereas $E-x^2/2\to\infty$. Hence this branch tends
uniformly to $y=-1$, covering every finite point of the lower arc.
At $y=0$ the two points $x=\pm\sqrt{2E}$ approach the horizontal
equatorial endpoints.

Fix an upper equatorial direction $(\cos\phi,\sin\phi)$ with
$0<\phi<\pi$ and $\cos\phi\ne0$. On its ray take
$(x,y)=c\epsilon^{-1/2}(\cos\phi,\sin\phi)$. Uniformly for $c$
in a fixed compact positive interval,
\[
 \epsilon\mathscr H_\epsilon(x,y)
     =\frac{c^2\cos^2\phi}{2}+o(1),
 \qquad 0<y<1/\epsilon.
\]
Values of $c$ on opposite sides of $1/|\cos\phi|$ bracket every
selected $E$. The intermediate value theorem supplies an intersection
on the bounded-well component, tending to that equatorial direction.
For the vertical direction, its upper turning point $y_+$ satisfies
$V_\epsilon(y_+)=E$ and
$y_+\ge E/(1+\epsilon)\to\infty$, since
$V_\epsilon(y)\le(1+\epsilon)y$ for $y\ge0$.
Every point of $\Gamma$ is therefore approached uniformly over the
selected energy band. Applying these pointwise estimates to a finite
net of the compact curve $\Gamma$ gives the reverse distance estimate
and proves \eqref{rw:lower:hausdorff}.

Choose $\epsilon_n\downarrow0$ so that the compact center families
lie in smaller neighborhoods whose closures lie in $\mathscr U$.
At each fixed $\epsilon_n$, choose one $\theta_n>0$ below the
root-persistence and path-positivity cutoffs for all three intervals,
small enough that the three perturbed cycles remain in $\mathscr U$,
and also $\theta_n<1/n$. Finite-time dependence over the compact
transverse center family lets us require, at the same time, that each
perturbed cycle have spherical Hausdorff distance less than $1/n$
from its corresponding center oval. Equation~\eqref{rw:lower:hausdorff}
then gives convergence of each of the three cycles to the entire
$\Gamma$.
In original coefficients this sequence is
\[
 \begin{gathered}
 B=0,\qquad m=\epsilon_n,\qquad
 a=\sqrt{\epsilon_n}\,h_{\epsilon_n}\theta_n,\\
 c=\sqrt{\epsilon_n}\,
   \bar c_{\epsilon_n}(h_{\epsilon_n}\theta_n,\theta_n,
                              \beta_{\epsilon_n}\theta_n),\qquad
 d=c+\sqrt{\epsilon_n}\,\beta_{\epsilon_n}\theta_n .
 \end{gathered}
\]
Reduce $\theta_n$ further so that $|\bar a|,|\bar c|<1$.
The primary scale remains exactly $\rho=\sqrt{\epsilon_n}$,
and all original coefficients tend to zero.
These fields lie in the original five-parameter family and hence
prove the lower bound in the full quadratic family, independently
of the upper bound and its forcing theorem.
\end{proof}

\begin{corollary}[Stability of the three cycles]\label{rw:lower:stability}
For the positive perturbation amplitudes in the construction of
Theorem~\ref{rw:three-simple-cycles}, the three selected cycles are,
in their nesting order, attracting, repelling and attracting.
\end{corollary}
\begin{proof}
At the three increasing simple roots of $P$, its derivative has
signs $+,-,+$. The two $C^1$ limits used above preserve these signs
for $F'(t)$, since $\theta,S_\epsilon,\delta_\epsilon>0$.
Equation~\eqref{rw:lower:actual-return-multiplier} and $U'>0$ then
give forward return multipliers less than one, greater than one,
and less than one. All are positive. Increasing $t$ moves inward
in the period annulus; reversing the order leaves this alternating
pattern unchanged.
\end{proof}
